\section{Literature}
\label{s:literature}

\subsection{\todo{Agglomeration and Marshall}}
\label{ss:lit-marshall}
It has been common to analyses of economic growth since \cite{Marshall1920}
to assume and qualitatively discuss some effect on firm-to-firm spillover of productivity.
A productivity spillover underpins well-established economic phenomena
such as the colocation of firms in urban clusters described in New Economic Geography literature
such as \cite{Krugman1991}.
It has been heavily analysed through industry and spatial density approaches
in spatial and urban econometric literature such as \cite{Glaeser1992}.

My investigation is rooted in \cite{Marshall1920}'s taxonomy,
which identifies three channels to productivity spillovers: (1) input-output costs, (2) labour pooling, and (3) knowledge spillovers.
There is no reason to think these are exhaustive. We might theorise that there may be competitive effects and cultural factors.
\cite{Krugman1991} also labels all these channels as 'centripetal' to an urban core; there are likewise 'centrifugal' factors that
push firms away from the urban core, such as destructive (zero-sum) competition, congestion, and saturation costs.
A further extension could develop models on these further channels and empirically test them.
Furthermore, these productivity spillovers are incredibly salient to thinking about production structures.
\par
In this investigation, I concentrate on (3) knowledge spillovers, tackling a difficult identification problem highlighted in Glaeser (2001).
This is more relevant to the knowledge-based services economy that dominates the UK.
Input-output costs and labour-pooling matter more for industrial manufacturing firms, especially the latter which was devised in response to
the phenomenon emergence of industry towns.
\par
Most agglomeration literature bundles all three channels together. Glaeser (2010) and Krugman (2010) identify how vital
identifying knowledge spillovers are. That need to estimate them well only grows with the modern, knowledge economy of the UK.
(Goldin, 2023) speculates that agglomeration effects have “increasing strength” attributable to the way in which
"knowledge workers in particular become more productive when located near to one another.
Information may be more available than ever, but the evidence suggests that it still tends to flow largely through in-person networks".

\subsection{\todo{New Economic Geography}}
\label{ss:lit-neg}

However, \cite{Krugman1991} prefers not to assume this effect
in their focus on the dynamic movement of firms.
Instead, heterogenous cluster policies, such as initial external differences in the number of workers
or firms are frequently assumed, with `forward and backward linkages' driving colocation
endogenously from there. This paper helpfully tests the presence of spillovers,
providing another centrifugal force directly incentivising firms to colocate. 

\subsection{\todo{Modern approaches}}
Specific approaches to investigating knowledge spillovers via innovation, instrumenting via patent publications comes from US evidence,
including (Jaffe, 1993), (Berkes and Ruben, 2021) and more recently \cite{Bloom2013}.
\par
Knowledge spillovers have been a well-evidenced factor of
managerial literature within the firm contributing to productivity gains.
A field-starting MIT study from \cite{Allen1977} observed an exponentially declining curve depicting the
frequency of communication between two individuals against their physical distance in an office.

\subsection{Peer effects econometrics}
\label{ss:lit-peer-effects-metrics}

The \cite{Manski1993} model is a common baseline in literature on peer effects.
It provides helpful structure to the study of joint behaviour of individuals
who are members of a group or network, and can be applied to any examination
of units with peers such as individuals, not merely firms.
\begin{align*}
    y_i \quad
        &=\quad c \quad
        + \underbrace{\alpha_g}         _{\substack{correlated \\\text{effect}}}
        + \quad\beta x_{i,g}\quad
        + \underbrace{\delta\bar y_{-i,g}}_{\substack{endogenous \\\text{effect}}}
        + \underbrace{\xi\bar x_{-i,g}}   _{\substack{exogenous  \\\text{effect}}}
        + \quad \varepsilon_i
\end{align*}
The model firstly decomposes spillovers to a direct behavioural, or \itx{endogenous} effect,
$y_{j}\to y_{i}$, where the outcome of interest among peers
affects that of the target individual.
Robust applications of the Manski model demonstrate credible approaches to measuring the \itx{endogenous effect}.
\cite{dustmann-damm2024:neighbourhoods} analyses movers to new neighbourhoods among individuals through
the `Danish Dispersal Policy' of refugee arrivals (1986-1998).
\cite{Chetty2016} measures a similar effect among high-poverty children through the US 1990s `Moving to Opportunity' scheme.
These studies exploit quasi-randomness in the location assignment of individuals and peers
for robust identification. As discussed in section \ref{ss:data-location-limits},
the dataset used for this analysis unfortunately offers no such granularity of features
to identify movers, let alone evaluating the exogeneity of their colocation.
The use of firm fixed effects and reasonable assumptions of location stickiness
are employed to compensate for this risk of location endogeneity.
\par
Additionally, the \cite{Manski1993} model requires contextual or \itx{exogenous} effects,
of which the variation of non-outcome of interest variables among peers
also directly affects that of the target individual, $x_j \to y_i$.
This brings sophistication to the nature of peer spillovers.
A simplified model only allowing endogenous effects through the variable of interest, $y_j \to y_i$
is nested within this structural model, but not identified due to the `Reflection Problem'
and as described in this analysis in section \ref{ss:static-identification}.
In this analysis' firm productivity model, the \itx{exogenous effect}
allows for an additional level of input factor capital and labour
to induce higher target TFP, better management practices or
more efficient combination of factor inputs, perhaps from observing
the mere high level of these factors in their peers,
without the peers themselves necessarily experiencing a productivity increase.
It is additionally a common specification requirement for every individual
to have an own-effect of exogenous characteristics, $x_i \to y_i$,
which corresponds to the standard production function model in this analysis of
$k_{i,t}, l_{i,t}\to y_{i,t}$.
\begin{align*}
    y_{i,g} &= c + \alpha_g + \beta x_{i,g} + \delta\bar y_{-i,g} + \xi\bar x_{-i,g} + \varepsilon_i
    \\
    \bar y_g &= c + \alpha_g + \beta\bar x_g + \delta\bar y_g + \xi\bar x_g + \bar\varepsilon_g
    \\
    (1-\delta)\bar y_g &= c + \alpha_g + (\beta+\xi)\bar x_g + \bar\varepsilon_g
    \\ \tag{large g assumption}
    \bar y_g &= \frac{c + \alpha_g}{1-\delta} +\frac{\beta+\xi}{1-\delta}\bar x_g+\frac{\bar\varepsilon_g}{1-\delta}
        \quad = \bar y_{-i,g}
    \\ \tag{substituting}
    y_{i,g} &= \frac {c + \alpha_g}{1-\delta} +\beta x_i + \frac{\xi+\delta\beta}{1-\delta}\bar x_{-i}
    + \frac\delta{1-\delta}\bar\varepsilon_g + \varepsilon_i 
    \\
    y_{i,g} &= \alpha_g + \pi_0 + \pi_1x_i + \pi_2 \bar x_{-i} + v_g
\end{align*}

The focus of \cite{Manski1993}'s linear-in-means (LMM) contribution is to highlight the `Reflection Problem',
an econometric trap.
The derivation give above shows that the initial structural equation
cannot be estimated directly with OLS due to the endogeneity of $\bar y_{-i,g}: \mathbb E[\bar y_{-i,g}\varepsilon_i]\neq 0$.
Transforming the equation to reduced form and omitting the endogenous regressor
leads to underidentification: 3 parameter estimates yielded for 4 structural parameters.
Causal variation in $y_{i,g}$ is unidentifiable as it is exactly collinear with all the other regressors.
This underidentification is a famous one, hindering analysis of peer dynamics including via spatial interaction.
\par
A key solution to this is the deployment of instruments, $z_{-i}\propto \bar y_{-i}\ \cap\ z_{-i} \not\propto y_{i}$.
An application employing credibly exogenous instruments for a peer effects measure in an individual crime setting
is given in \cite{dustmann-landerso2024:gangs-babyfathers}.
This study uses the sex-identity at birth of newborns as a valid instrument
for the criminal behaviour of their babyfather gang-leaders,
who serve as an influential peer on the target members of their gangs respectively.
The equivalent valid instrument in this analysis which is correlated with the peer effect
but has no direct effect on the target outcome of interest
is an additional spatial lag of \itx{exogenous} regressors in the static model.
This means the capital and labour of a firm $k$ who is a `peers of peers': a peer of firm $j$,
but not of firm $i$, the last of which is the target firm of interest.
In this analysis' dynamic model, deep temporal lags $y_{-i,t-2}\propto y_{-i,t}\ \cap\ y_{-i,t-2}\not\propto y_{i, t}$
are also as a valid instrument under an assumption of weak exogeneity of predetermined regressors.
\par
A modern paper, \cite{Bramoulle2009}, explores different identification techniques and their rules for validity
for group and network structures. These structures are more complex than the \cite{Manski1993} model,
which by using a variable of means $\bar y_{-i,g}$, implicitly assumes simple, `leave-one-in' group averages
as the peer \itx{endogenous} regressor.
\cite{Moffitt2001} makes a modification to drop the `large $g$ assumption'
and instead exclude the target observation $y_{i,g}$ from calculation of the group mean $y_g$,
which provides a helpful conceptual move from estimating a `neighbourhood effect',
for which the target is still a contributor,
towards more of a collection of influential peers.
Nevertheless, this is still insufficient to offer identification.
\par
The paper subsequently explores many complex interaction structures summarised by the weighting matrix $\mathbf W$,
the $n$-dimensional interactions matrix between firms.
I discuss and illustrate structures of $\mathbf W$ in detail in section \ref{ss:data-w-groups}.
The structures of $\mathbf W$ are key to identification.
The contributions from this paper are, in short, that complex $\mathbf W$ matrices which vary group sizes ($s_g$)
and network diameters ($s_d$) are key to discovering valid spatial lag instruments.
Fortunately, these structures resemble realistic interactions common to empirical data,
and are applicable to this analysis' pattern and layout of firm interactions.
We should note in application that this strategy
continues to assume exogenous assignment of spatial assignments
and so target-peer interactions.

\wraptable
    {Spatial econometric identification rules for weighting matrix $\mathbf W$}
    {./figures/2_tab_bdf_identification}
    {
        Adapted from \citepalt{Bramoulle2009} with notation
        revised and abridgments to improve relevance to this analysis.
        (a) \citepalt{Manski1993}, in this model $\mathbf W^2=\mathbf W$ for linear dependence.
        (b) \citepalt{Moffitt2001}, in this model, $\mathbf W^2 = \frac1{s-1}\mathbb I + \frac{s-2}{s-1}\mathbf W$ for linear dependence.
        (c) \citepalt{Lee2007}.
        Structures under a \textbf{Proposition} header must follow the \textbf{Proposition}'s required conditions.
        $x$ denotes an \itx{exogenous} regressor, in this firm analysis this means $k_{i,t}$ or $l_{i,t}$.
        It is comparable to the vector of regressors $\mathbf x_{i,t}$ that I use to denote
        the vector of both of them.
    }

The identification rules from \cite{Bramoulle2009} are summarised in \tablerefp{tab:bdf_identification}.
Critically, once certain necessary conditions for the group and network structures are met,
an additional spatial lag on the \itx{exogenous} regressors, $\mathbf W^2x$,
tends to be a valid instrument.
This analysis uses \tablerefp{tab:bdf_identification}'s approaches to instruments
to strip out the endogeneity of a spatial peer effects model.